%! Author = murfe
%! Date = 11.09.2026

\chapter{Diskussion und Fazit}
\label{ch:diskussion}
Dieses Kapitel ordnet die Ergebnisse in den Gesamtkontext der Arbeit ein.
Dazu werden die Forschungsfragen beantwortet und die Grenzen der Ausarbeitung diskutiert.
Abschließend wird ein Ausblick über weiterführende Untersuchungen und Verbesserungen des Prototyps gegeben.

\section{Gewonnene Erkenntnisse}
\label{sec:erkenntnisse}
Aus der Entwicklung und Evaluation des Systems ergeben sich mehrere zentrale Erkenntnisse für die Gestaltung fehlertoleranter
Hintergrundverarbeitung.
Dieses Kapitel beantwortet in diesem Zusammenhang die Forschungsfragen und führt zusätzlich gewonnene Erkenntnisse auf.

\subsection{Beantwortung der Hauptforschungsfrage}
\label{subsec:beantwortungHauptfrage}
Auf Grundlage des entworfenen und evaluierten Prototyps lässt sich die Hauptforschungsfrage für das betrachtete Fehlermodell dahingehend
beantworten, dass zur Erfüllung mehrere aufeinander abgestimmte Mechanismen erforderlich sind.
Das persistente Zustandsmodell bildet dafür die Grundlage und ist besonders relevant für die konkurrierende Verarbeitung,
Recovery und Konsistenz der Daten.
Transaktionen und zeilenbasierte Sperren synchronisieren konkurrierende Zugriffe und verhindern das Zurückbleiben fehlerhafter
Teildaten im Fehlerfall.
Idempotenz wird an zwei Stellen benötigt: Zum einen wird mithilfe von \gls{uuid}s das mehrfache Anlegen identischer Jobs verhindert.
Zum anderen werden eindeutige Ergebnisse durch das Koordinieren von Zugriffen auf Basis von Fencing erzielt.
Während Retry-Strategien die Wiederholung transient fehlgeschlagener Operationen ermöglichen, sorgt
die lease-basierte Recovery mit regelmäßiger Abfrage des Jobbestands dafür, dass nach einem Worker-Ausfall verwaiste
Jobs wieder für die Verarbeitung freigegeben werden können.

\subsection{Beantwortung der Teilfragen}
\label{subsec:beantwortungTeilfragen}
Für \ref{item:tf1} wurden in \cref{sec:zusicherungen} in Abhängigkeit vom betrachteten Fehlermodell konkrete Zusicherungen abgeleitet.
Die Evaluation zeigt, dass die Teilausfälle des betrachteten Fehlermodells durch die Umsetzung dieser Zusicherungen behandelt werden können.
Insbesondere gilt das für Infrastrukturausfälle wie Netzwerkfehler oder Worker-Crashes.

Die Untersuchung zeigt bezüglich \ref{item:tf2}, dass diese Mechanismen unterschiedliche Fehlerklassen adressieren und sich daher gegenseitig ergänzen.
Retries ermöglichen die automatische Neuausführung von Operationen, wenn ein Fehler als transient eingestuft wird und die
ausführende Komponente weiterhin aktiv ist.
Recovery übernimmt dagegen die Wiederaufnahme von Jobs, bei denen kein lokaler Retry mehr möglich ist, beispielsweise durch
den Absturz eines aktiven Prozesses.
Idempotenz ergänzt diese beiden Fehlertoleranzstrategien, indem sie sicherstellt, dass die erneute Ausführung einer Operation
keine unerwünschte Mehrfachwirkung erzeugt.

In Hinsicht auf \ref{item:tf3} zeigen die Ergebnisse, dass kurze, durch exklusive Sperren geschützte Transaktionen zur Zustandsänderung
ausreichen, um die entscheidenden konkurrierenden Operationen auf Datenbankebene zu synchronisieren, sofern ein passendes
Zustandsmodell für den Lebenszyklus der Jobs verwaltet wird.
Beim Zustandsübergang eines Jobs verhindert die Kombination aus Transaktion und Zeilensperre, dass mehrere Instanzen denselben
Job gleichzeitig auf Datenbankebene bearbeiten können.
Wird vor dem Abschluss eines Jobs zusätzlich sein aktueller Zustand und seine Besitzverhältnisse geprüft, können konkurrierende
Operationen ohne das Auftreten von Inkonsistenzen synchronisiert werden, obwohl kurze Transaktionen verwendet werden.

Für \ref{item:tf4} zeigt die Untersuchung, dass hierfür die Kombination aus persistentem Zustandsmodell, zeitbasiertem Lease
und Recovery-Mechanismus wesentlich ist.
Fällt ein Worker während der Verarbeitung aus, verbleibt der Job im persistenten Job-Store zunächst im Zustand \textit{Running}.
Nach Ablauf des Leases kann der Recovery-Service erkennen, dass die Verarbeitung nicht innerhalb des vorgesehenen Zeitraums
abgeschlossen wurde und den Job erneut zur Verarbeitung freigeben.
Bereits erfolgreich persistierte Jobdaten bleiben aufgrund der \textit{Durability}-Eigenschaft der relationalen Datenbank auch über
einen Ausfall des Workers hinweg erhalten, sodass ein ausgefallener Worker nicht zum Verlust des persistent gespeicherten Jobs führt.

\subsection{Weitere architektonische Erkenntnisse}
\label{subsec:weitereErkenntnisse}
Eine zentrale Erkenntnis der Untersuchung ist, dass eine relationale Datenbank als Job-Store nicht ausschließlich als
persistenter Speicher, sondern zugleich als Koordinationsmechanismus für die nebenläufige Verarbeitung und für Teile der
Konsistenzsicherung der Jobs und Ergebnisse eingesetzt werden kann.
Dadurch lässt sich ein wesentlicher Teil der geforderten Zusicherungen bereits durch die Datenbank realisieren.

Für die untersuchte Systemgröße ergibt sich daraus als Antwort auf die in der Einleitung aufgeworfene Frage ein weiterer Vorteil:
Die Architektur benötigt keinen zusätzlichen Message Broker als eigenständige Infrastrukturkomponente.
Insbesondere entfällt dadurch architektonischer Aufwand zur Konsistenzerhaltung zwischen Datenbank und Broker.
Die Nutzung einer gemeinsamen relationalen Datenbank ermöglicht es stattdessen, zusammengehörige Änderungen über lokale
\gls{acid}-Transaktionen abzusichern.
Dies reduziert die infrastrukturelle Komplexität, hat aber auch Auswirkungen auf die Skalierbarkeit, die im folgenden Kapitel
thematisiert werden.

\section{Einschränkungen und Grenzen der Arbeit}
\label{sec:einschraenkungen}
\textbf{Prototypischer Charakter und begrenztes Fehlermodell.}
Das entwickelte System dient der Untersuchung ausgewählter Fehlertoleranzmechanismen und stellt kein vollständig
produktionsreifes Job-Processing-System dar.
Dementsprechend besteht Optimierungsbedarf bei der Parametrisierung, beispielsweise bezüglich der maximalen Anzahl an Wiederholungsversuchen,
der initialen Verzögerung oder des Backoff-Faktors.
Besonders die Konfiguration der Retry-Policies bezüglich der Klassifizierung transienter Fehler ist als \textit{Proof of concept}
und nicht als vollständige Lösung zu verstehen.

Bezüglich \ref{zus:verarbeitungsgarantie} bedeutet die At-Least-Once-Verarbeitungsgarantie eines Jobs nicht, dass für jeden Job zwangsläufig
ein Ergebnis erzeugt wird.
Vielmehr bedeutet das, dass die Verarbeitung nach endlicher Zeit abgeschlossen wird, aber nicht zwangsweise erfolgreich.
Bei wiederholtem Auftreten eines Fehlers wird die Verarbeitung nach der festgelegten maximalen Anzahl von Versuchen bewusst
beendet und der Job in einen definierten Fehlerzustand überführt.
Damit wird eine unbegrenzte Wiederholung verhindert, zugleich kann ein Job trotz seiner dauerhaften Persistierung letztlich als fehlgeschlagen enden.
Die Zusicherung bezieht sich somit auf die Vermeidung eines unbeabsichtigten Verlusts von Jobs und nicht auf eine Garantie
des erfolgreichen Abschlusses jeder Verarbeitung.

Des Weiteren werden Fehler außerhalb des definierten Modells nicht betrachtet, da Fehlertoleranz nicht als absolute Eigenschaft
eines Systems verstanden werden kann \parencite[Kap.~1 Abschnitt \glqq Reliability\grqq]{kleppmann_designing_2017}.
Die implementierten Mechanismen adressieren gezielt die im Rahmen dieser Arbeit betrachteten Fehlerfälle.
Sie reduzieren die Auswirkungen dieser Fehler, können jedoch keine allgemeingültige Garantie gegenüber jedem denkbaren Ausfall liefern.
Die Aussagekraft der erzielten Ergebnisse ist daher fest an das definierte Fehlermodell gebunden.

\textbf{Zeitbasierte Leases ohne Heartbeat.}
Die prototypische Implementierung verwendet zeitbasierte Leases und verzichtet auf einen Heartbeat-Mechanismus.
Dadurch kann ein ungewöhnlich aufwendiger Job oder ein aktiver, aber langsamer Worker dazu führen, dass die Verarbeitung
eines Jobs immer wieder durch die Recovery abgebrochen und niemals erfolgreich beendet wird.
Ein Heartbeat-Mechanismus könnte diesem Problem entgegenwirken, indem der Worker die Gültigkeit seines Leases regelmäßig verlängert.
Dadurch ließe sich das Risiko reduzieren, dass lang laufende Verarbeitungen fälschlicherweise als ausgefallen eingestuft werden.
Dem steht jedoch eine zusätzliche Belastung des Job-Stores gegenüber, da während der Verarbeitung regelmäßig Schreibzugriffe zur
Aktualisierung des Leases erforderlich wären.
Die Wahl zwischen diesen beiden Lösungen stellt also eine Abwägung zwischen einer geringeren Belastung des Job-Stores und einer
höheren Robustheit gegenüber ungewöhnlich langen Verarbeitungszeiten dar.
Bei der Nutzung von Leases ohne Heartbeat muss außerdem die Wahl der Lease-Dauer beachtet werden.
Eine kurze Lease-Dauer ermöglicht eine schnellere Recovery, erhöht jedoch das Risiko, dass noch aktive Verarbeitungen vorzeitig
zur erneuten Verarbeitung freigegeben werden.

\textbf{Begrenzte Aussagekraft der empirischen Evaluation.}
Die Evaluation erfolgt in einer kontrollierten Testumgebung und basiert auf einer endlichen Anzahl von Testwiederholungen.
Insbesondere für nebenläufigkeitsabhängige Szenarien kann daher nicht ausgeschlossen werden, dass unter anderen Ausführungsbedingungen
weitere bisher nicht beobachtete Effekte auftreten.
Die wiederholte Ausführung der entsprechenden Tests erhöht die Wahrscheinlichkeit, relevante Nebenläufigkeitseffekte zu erfassen,
stellt jedoch keinen Beweis für deren vollständiges Ausbleiben dar.
Des Weiteren kann mit Tests im Allgemeinen nur das Vorhandensein eines Fehlers bewiesen werden, nicht jedoch die Abwesenheit.
Trotz der erfolgreichen Tests kann die Funktionsfähigkeit des Systems daher auch bezüglich des betrachteten Fehlermodells
nicht vollständig garantiert werden.
Ein vollständiger Ausschluss solcher Effekte würde eine formale Verifikation erfordern.

\textbf{Keine Untersuchung von Performance und Skalierbarkeit.}
Die Arbeit untersucht die Funktionsfähigkeit der Fehlertoleranzmechanismen und die Konsistenz des entwickelten Systems,
nicht dessen Verhalten unter hoher Last.
Die Verwendung eines datenbankbasierten Job-Stores mit Sekundärspeicher kann bei einer hohen Anzahl an Zugriffen aufgrund
der Latenz und für den parallelen Zugriff benötigten Ressourcen zum limitierenden Faktor werden
\parencite[Kap.~10 Abschnitt \glqq Example: analysis of user activity events\grqq]{kleppmann_designing_2017}.
Für größere Systeme kann der Einsatz von Message Brokern hinsichtlich Durchsatz und Skalierbarkeit deshalb Vorteile bieten.
Daher wäre ein Performancevergleich des entwickelten Prototyps mit einer Lösung, die auf einem Message Broker basiert, besonders
interessant.
Da die implementierten Mechanismen primär die zuverlässige Verarbeitung bereits persistierter Jobs adressieren und nicht
die vorgelagerte Lastverteilung bei stark erhöhtem Anfrageaufkommen, ist auch das Verhalten bei hohen Eingangslasten spannend.
Für entsprechende Szenarien werden in der Praxis häufig zusätzliche Mechanismen zur Lastbegrenzung eingesetzt \parencite{microsoft_corporation_hrsg_queue-based_2026}.

Die verwendeten Retry-Strategien nutzen außerdem ein exponentielles Backoff, jedoch keine zusätzliche Randomisierung der Wartezeiten.
Da sich die Last dadurch gleichmäßig auf der Zeitachse nach rechts verschiebt, können neue Anfragen, die mitten in einem
Schwall an Retries ankommen, zur Überlastung des Systems führen \parencite[Kap.~3 Abschnitt \glqq Jitter\grqq]{newman_building_2026}.
Eine Erweiterung um ein Jitter könnte das Problem besonders bei einer größeren Anzahl gleichzeitig betroffener Anfragen weiter reduzieren.

\section{Fazit}
\label{sec:fazit}
Ziel dieser Arbeit war es, zu untersuchen, welche Architekturmechanismen und Entwurfsentscheidungen es ermöglichen, in einem
Job-Processing-Backend definierte Konsistenzzusicherungen auch unter typischen Teilausfallszenarien einzuhalten.
Die Untersuchung zeigt, dass sich diese Zusicherungen im betrachteten Fehlermodell durch das Zusammenspiel eines persistenten
Zustandsmodells, transaktionsgesicherter Datenbankoperationen mit geeigneten Locking-Mechanismen, Idempotenz, Retry-Strategien
sowie einer lease-basierten Recovery mit Besitz- und Lease-Prüfung allein auf Basis einer relationalen Datenbank umsetzen lassen,
ohne eine zusätzliche Broker-Infrastruktur zu benötigen.
Die formulierten Teilfragen konnten durch die hergeleiteten Zusicherungen, den darauf aufbauenden Architekturentwurf und
die durchgeführte Evaluation beantwortet werden.
Außerdem stützen die in der Evaluation durchgeführten Tests die formulierten Hypothesen.

\section{Ausblick}
\label{sec:ausblick}
Aus den Ergebnissen und den aufgezeigten Einschränkungen ergeben sich mehrere Ansatzpunkte für weiterführende Untersuchungen.

Ein zentraler nächster Schritt wäre eine Untersuchung der Performance und Skalierbarkeit des Systems.
Die vorliegende Evaluation konzentriert sich auf die Korrektheit und Fehlertoleranz und betrachtet keine hohe Systemlast.
Insbesondere wäre zu untersuchen, wie sich eine steigende Anzahl gleichzeitig eintreffender und zu verarbeitender Jobs auf Datenbanklast,
Latenz und Durchsatz auswirkt und ab welchem Systemumfang eine Message-Broker-Lösung Vorteile gegenüber dem
datenbankbasierten Ansatz bietet.

Darüber hinaus könnte die Retry- und Recovery-Strategie weiterentwickelt werden.
Die Ergänzung eines Jitters zur exponentiellen Backoff-Strategie könnte eine weitere Reduktion synchronisierter Wiederholungsversuche ermöglichen.
Ebenso wäre ein empirischer Vergleich der verwendeten Lease-Strategie mit einem Heartbeat-basierten Ansatz interessant,
insbesondere im Hinblick auf die Unterscheidung zwischen tatsächlich ausgefallenen und lediglich langsam arbeitenden Workern.

Ein weiterer Ansatzpunkt zur Weiterentwicklung des Systems ist der Umgang mit externen Seiteneffekten.
Während die vorliegende Arbeit die konsistente Persistierung von Job-Ergebnissen betrachtet, können bei einer Wiederaufnahme
eines Jobs nicht-idempotente externe Operationen, beispielsweise der Versand einer E-Mail, mehrfach ausgeführt werden.
Da solche Operationen nicht gemeinsam mit der lokalen Datenbanktransaktion atomar ausgeführt werden können, wäre zu untersuchen,
wie die Absicherung gegen die Mehrfachausführung solcher Seiteneffekte in das bestehende System integriert werden kann.
